\chapter{Menerima Pengaruh: Power Sharing dalam Hubungan}

\begin{insightbox}
\textit{``Hubungan yang sehat bukanlah tentang siapa yang benar, tapi tentang bagaimana kita saling menghargai perspektif. Menerima pengaruh adalah fondasi dari kekuasaan bersama.''} — Dr. John Gottman
\end{insightbox}

\section{Apa itu Menerima Pengaruh?}

\keyterm{Menerima pengaruh} adalah kemampuan untuk membuka diri terhadap perspektif, perasaan, dan kebutuhan pasangan---terutama saat berkonflik. Ini bukan berarti selalu mengalah atau kehilangan diri sendiri, melainkan memiliki sikap keterbukaan: \textit{``Mungkin kamu punya poin yang valid. Mari kita diskusikan.''}

\subsection{Kenapa Ini Penting?}

Penelitian Gottman yang melibatkan 130 pasangan baru menikah menemukan:

\begin{praktikbox}[Fakta Ilmiah]
\begin{itemize}
    \item Pria yang \textbf{menerima pengaruh} dari istrinya memiliki perkawinan yang \textbf{signifikan lebih bahagia}
    \item Pria yang \textbf{menolak pengaruh} istri memiliki \textbf{81\% kemungkinan lebih tinggi} untuk bercerai
    \item Kesenjangan gender ini terkait dengan \textbf{socialisasi tradisional}, bukan biologi
\end{itemize}
\end{praktikbox}

Ini bukan berarti wanita selalu benar atau pria selalu salah. Ini tentang \textbf{power dynamics}---bagaimana tradisi patriarkal sering mengajarkan pria untuk ``memimpin'' tanpa mendengarkan, yang justru merusak hubungan.

\section{Karateristik Menerima vs Menolak Pengaruh}

\begin{center}
\renewcommand{\arraystretch}{1.6}
\begin{tabular}{|>{\raggedright}p{6cm}|>{\raggedright\arraybackslash}p{7cm}|}
\hline
\rowcolor{primary!20}
\textbf{MENOLAK PENGARUH (Toxic)} & \textbf{MENERIMA PENGARUH (Healthy)} \\
\hline
``Aku tahu apa yang terbaik'' & ``Ceritakan pemikiranmu'' \\
\hline
Memotong pembicaraan pasangan & Mendengarkan sampai selesai \\
\hline
Meyakinkan pasangan salah & Mencari validity dalam argumen pasangan \\
\hline
Body language defensive (mencibir, mata guling) & Body language terbuka (kontak mata, angguk) \\
\hline
Mengabaikan perasaan pasangan & Validasi: ``Aku mengerti kamu merasa...'' \\
\hline
Kompetisi: harus menang & Kolaborasi: mencari win-win \\
\hline
Menyalahkan pasangan & Mengambil tanggung jawab pribadi \\
\hline
Menggunakan kekerasan/emosi untuk kontrol & Menggunakan diskusi untuk solusi \\
\hline
\end{tabular}
\end{center}

\section{Self-Assessment: Seberapa Terbuka Kamu?}

\begin{exercisebox}[Test: Tingkat Keterbukaan Terhadap Pengaruh]
\textbf{Petunjuk:} Jawab setiap pertanyaan dengan skala 1-5 (1 = Tidak pernah, 5 = Selalu)

\begin{enumerate}[leftmargin=*]
    \item Saat pasangan mengajukan pendapat berbeda, aku benar-benar mendengarkan tanpa memotong.
    \item Aku bisa mengubah pendapatku jika pasangan memberikan argumen yang kuat.
    \item Saat konflik, aku bertanya ``Apa yang kamu butuhkan?'' sebelum memaksakan keinginanku.
    \item Aku mengakui kesalahanku dan meminta maaf dengan tulus.
    \item Aku menghormati perasaan pasangan meski logika mengatakan sebaliknya.
    \item Aku tidak merasa ``kalah'' saat mengikuti saran pasangan.
    \item Aku mencari middle ground daripada insist pada solusiku sendiri.
    \item Saat pasangan mengkritik, aku mencari kebenarannya bukannya defensive.
    \item Aku melibatkan pasangan dalam keputusan besar (finansial, karir, dll).
    \item Aku bisa bilang ``Kamu benar'' tanpa merasa terancam.
    \item Saat pasangan marah, aku tetap tenang dan mencari tahu akar masalahnya.
    \item Aku menghargai keahlian pasangan di bidang yang bukan kekuatanku.
\end{enumerate}

\textbf{Scoring:}
\begin{itemize}
    \item \textbf{48-60:} Excellent --- Kamu ahli dalam power sharing
    \item \textbf{36-47:} Good --- Ada ruang untuk improvement
    \item \textbf{24-35:} Fair --- Perlu kerja keras dalam area ini
    \item \textbf{12-23:} Needs Work --- Risiko hubungan tinggi tanpa perubahan
\end{itemize}
\end{exercisebox}

\section{Mengapa Pria Cenderung Menolak Pengaruh?}

\begin{warningbox}[Bukan Salah Pria Secara Individual]
Ini adalah \textbf{socialization issue}---bagaimana budaya membentuk ekspektasi gender, bukan kegagalan pribadi.
\end{warningbox}

\subsection{Faktor Budaya dan Sosial}

\begin{enumerate}[leftmargin=*]
    \item \textbf{Socialization Masculine:} Sejak kecil, laki-laki diajarkan untuk ``kuat'', ``tidak menyerah'', ``memimpin''---sering diartikan sebagai tidak boleh dipengaruhi.
    
    \item \textbf{Power = Identity:} Bagi banyak pria, kemampuan mengendalikan situasi terkait dengan self-worth. Menerima pengaruh terasa seperti kehilangan identitas.
    
    \item \textbf{Fear of Appeasing:} Takut dianggap `` whipped'' atau ``tidak berprinsip'' jika mengikuti pasangan.
    
    \item \textbf{Logical Processing:} Pria cenderung lebih fokus pada logika dan solusi, mengabaikan emotional content yang sering jadi inti argument wanita.
    
    \item \textbf{Conflict as Competition:} Olahraga/bidang kompetitif mengajarkan ``menang adalah segalanya'', sehingga konflik dianggap pertandingan.
\end{enumerate}

\subsection{Paradoksnya}

Paradoksnya: \textbf{Pria yang menerima pengaruh justru dihormati lebih} oleh pasangannya. Mereka dilihat sebagai:\
$\bullet$ Lebih dewasa emotionally\\
$\bullet$ Lebih confident (karena secure people don't need to control)\\
$\bullet$ Lebih menarik secara romantis\\
$\bullet$ Role model yang lebih baik untuk anak

\section{Cara Menerima Pengaruh: Step-by-Step}

\begin{praktikbox}[Framework: DEAR untuk Menerima Pengaruh]
\begin{enumerate}[leftmargin=*]
    \item \textbf{D}elay --- Tahan diri 6 detik sebelum merespons
    \item \textbf{E}xplore --- Tanyakan: ``Ceritakan lebih lanjut'' atau ``Apa yang kamu butuhkan?''
    \item \textbf{A}cknowledge --- Validasi: ``Aku mengerti kamu merasa...''
    \item \textbf{R}espond --- Berikan respon yang menunjukkan kamu mempertimbangkan
\end{enumerate}
\end{praktikbox}

\subsection{Frase-frase yang Menerima Pengaruh}

\begin{tipbox}[Kalimat Pembuka Keterbukaan]
\begin{itemize}
    \item ``Itu perspektif yang belum aku pikirkan.''
    \item ``Kamu punya poin yang valid.''
    \item ``Aku salah. Kamu benar.''
    \item ``Bagaimana kalau kita coba cara kamu?''
    \item ``Aku lihat kenapa kamu merasa begitu.''
    \item ``Mari cari solusi yang works untuk kita berdua.''
    \item ``Apa yang bisa aku lakukan untuk membantu?''
    \item ``Ceritakan lebih banyak tentang apa yang kamu rasakan.''
\end{itemize}
\end{tipbox}

\section{Attachment Styles dan Menerima Pengaruh}

\begin{center}
\renewcommand{\arraystretch}{1.5}
\begin{tabular}{|>{\raggedright}p{2.5cm}|>{\raggedright}p{5cm}|>{\raggedright\arraybackslash}p{5.5cm}|}
\hline
\rowcolor{primary!20}
\textbf{Attachment} & \textbf{Biasanya} & \textbf{Strategi Mengatasi} \\
\hline
\textbf{Anchor} & Naturally good at accepting influence; balance give-take & Lanjutkan, bantu pasangan yang struggle \\
\hline
\textbf{Island} & Resist influence (fear of losing autonomy); logical defense & Latihan: Acknowledge feelings first, then discuss logic \\
\hline
\textbf{Wave} & May accept too much influence (people-pleasing) then resent & Set boundaries; practice saying ``I need'' not just ``Okay'' \\
\hline
\end{tabular}
\end{center}

\section{Kapan Menolak Pengaruh adalah Wajar?}

\begin{warningbox}[Bukan Selalu Mengalah]
Menerima pengaruh \textbf{bukan} berarti:
\begin{itemize}
    \item Mengikuti segala permintaan pasangan tanpa batas
    \item Mengorbankan values atau principles fundamental
    \item Toleransi terhadap abuse atau manipulasi
    \item Kehilangan identitas diri
\end{itemize}
\end{warningbox}

\subsection{Menolak Pengaruh Secara Sehat}

Ada cara menolak yang tetap menghargai:
\begin{enumerate}[leftmargin=*]
    \item \textbf{Acknowledge dulu:} ``Aku mengerti kenapa kamu ingin ini.''
    \item \textbf{Jelaskan batasan:} ``Tapi ada sesuatu yang tidak bisa aku kompromikan.''
    \item \textbf{Tawarkan alternatif:} ``Bagaimana kalau...''
\end{enumerate}

\begin{exercisebox}[Latihan: Accepting Influence Practice]
\textbf{Skenario:} Pasanganmu ingin menghabiskan liburan dengan keluarganya, sedangkan kamu ingin liburan berdua.

\textbf{Respon yang Menolak Pengaruh:}
\begin{itemize}
    \item[$\square$] ``Keluargamu memang selalu lebih penting.''
    \item[$\square$] ``Aku sudah bilang berdua, jadi berdua.''
    \item[$\square$] ``Kamu tidak pernah menghargai keinginanku.''
\end{itemize}

\textbf{Respon yang Menerima Pengaruh:}
\begin{itemize}
    \item[$\square$] ``Aku mengerti kamu kangen keluarga.''
    \item[$\square$] ``Bagaimana kalau kita bagi: 3 hari keluarga, 4 hari berdua?''
    \item[$\square$] ``Apa yang membuat liburan keluarga penting bagimu tahun ini?''
\end{itemize}

\textbf{Latihan:} Diskusikan dengan pasangan: \textit{``Apa satu area di mana kamu merasa sulit menerima pengaruhku?''}
\end{exercisebox}

\section{Ringkasan Actionable}

\begin{tipbox}[Daily Practices]
\begin{enumerate}
    \item \textbf{Morning Check:} Sebelum memutuskan sesuatu, tanyakan pendapat pasangan.
    \item \textbf{Active Listening:} Saat pasangan bicara, tahan diri 6 detik sebelum merespons.
    \item \textbf{Acknowledgment Log:} Setiap hari, akui satu hal yang pasangan benar tentang.
    \item \textbf{Weekly Review:} Tanyakan ``Apa yang bisa aku lakukan lebih baik dalam mendengarkanmu?''
\end{enumerate}
\end{tipbox}
